%----------------------------------------------------------------------------------------
% LaTeX 
% TeXstudio 2.12.22 (git 2.12.22)
% author: Yuan Zhang
% yuan.zhang119@imperial.ac.uk
%----------------------------------------------------------------------------------------
\documentclass[11pt]{article}
\usepackage[utf8]{inputenc}
\usepackage{geometry}
\geometry{a4paper, portrait, margin=3cm}
\usepackage{graphicx}
\usepackage{amsmath, amssymb}
\usepackage{lineno}
\usepackage{cite}

\usepackage{float} % set figure location

\usepackage{helvet}
\renewcommand{\familydefault}{\sfdefault}
\usepackage{amsmath}
\usepackage{setspace}
\usepackage{indentfirst}
\usepackage{siunitx}
\usepackage{parskip}
\usepackage[labelfont=bf]{caption}
\usepackage[authoryear]{natbib}
\usepackage[T1]{fontenc}
\usepackage{multirow}
\usepackage{multicol}
\usepackage{rotating}



\linespread{1.5}

\begin{document}
	
	\begin{titlepage}
		
		\begin{center}
			
			\begin{figure}[t]
				\centering
				\includegraphics[width = \textwidth]{../Materials/logo.jpg}
			\end{figure}
			
			\vspace{1cm}
			\Huge
			\textbf{The Effect of Physiological Mismatches on The Temperature-Dependent Fitness of Disease Vectors }\\
			
			\vspace{3cm}
			\Large
			
			\textbf{Yuan Zhang}\\   
			\textbf{Auguest 2020}\\
			
			
			
			
			\textbf{A thesis submitted for the partial fulfillment of the requirements for the degree of Master of Science at Imperial College London}\\
				
				
			\textbf{Formatted in the journal style of Ecology Letters.
}\\
			\textbf{Submitted for the MSc  in Computational Methods in Ecology and Evolution}\\
			\vspace*{2cm}	

			
			
		\end{center}
		
	\end{titlepage}
	
	\clearpage
	\tableofcontents
	
	
	\linenumbers %show line numbers
	\doublespacing
	

\newpage

\section*{Abstract}

\section{Introduction}  
Corona Virus Disease 2019 (COVID-19), a new type of Coronavitus, accompanied as human-to-human transmission,  has became as a serious public health threat. Although similar to SARS-CoV and MERS-CoV, the COVID-19 is quite different \citep{R6, R7, R8}. The rapidly increasing evidence of human-to-human transmission suggest that the virus is more contagious than SARS-CoV and MERS-CoV \citep{R11}. Based on statistic by WHO(2020), the number of confirmed cases worldwide has exceeded 200 000. Besides, it also can be deadly for massive alveolar damage and progressive respiratory failure \citep{R11}. Although it has been found that 96\%  COVID-19 matched at whole genome level to a bat coronavirus\citep{R14}, other aspects such as transmission methods and affecting factors of this Coronavitus are unknown. In short, it is essensial to study this Coronavitus by suitable modeling.

\paragraph{} Environmental factors  (Temperature, humidity) have proven to be important influencing factors on epidemic disease \citep{R1}.  Some literature indicates that a high temperature and high humidity climate can reduce virus transmission  \citep{R3} . However,in fact, there are still many confirmed diagnoses in tropical countries such as Singapore and Malaysia. Therefore, the influence of environmental factors is complex and unclear. Few literatures has invesitigate this aspect and one fit empirical data seems not well  \citep{R1}.

\paragraph{} This study aims to investigate roles of Humidity and Temperature in COVID-19 infection dynamics. The proposed questions are:

(1) Are Coronavirus infections phenomenological data in most cities showing similar trends in global scale? 

(2) Whether environmental factors (focus on Temperature, humidity) have correlations with the number of diagnoses ? 

(3) If the temperature and humidity factors are added to the existing model to modified the model, will it be more suitable for a large number of sample data?

\section{Proposed Methods}
\subsection{Data and preparation}
Infections in different countries and regions;
Different temperature and humidity around the world
\subsection{ Analysis}
\subsubsection{ Environmental  impact}
First analyze the correlation between environmental factors and the number of infected people to see if you can find a rule from them(focus on the inflection point, minimum infection corresponding temperature and humidity). Some former researches has considered SARS and Environmental factors \citep{R9, R10}.

\subsubsection{ Fitting model and Modified model}
Fitting empirical data based on existing models:

(1) Weibull distribution using the Maximum Likelihood Estimation (MLE) method \citep{R2} 
This is linear regression, not true and the fit is macroscopically bad.

(2) Microsimulation model to two countries: the UK  and the US \citep{R4}

(3) A model based on aggregation of individuals according to disease status \citep{R15}

Try to change the existing model through adding environmental factors, and then look at the degree of fit. 

\section{Anticipated outputs and outcomes}
The improved model with the environmental factor model is more suitable for global data. 

\section{Project feasibility supported by a timeline of tasks}
\begin{figure}[thpb]
	\centering


\end{figure}
\section{The itemized budget}
Noun
\newpage

\bibliographystyle{agsm}
\bibliography{proposal}

\end{document}
